\documentclass[11pt,a4paper]{article}
\usepackage[T1]{fontenc}
\usepackage[utf8]{inputenc}
\usepackage{lmodern,microtype}
\usepackage[margin=25mm,headheight=15pt]{geometry}
\usepackage{amsmath,amssymb,amsthm,mathtools}
\usepackage{booktabs,longtable,array,enumitem}
\usepackage{xcolor,listings,xurl,hyperref,fancyhdr,needspace}
\definecolor{ink}{HTML}{24364B}
\definecolor{link}{HTML}{245B7B}
\hypersetup{colorlinks=true,linkcolor=link,urlcolor=link,citecolor=link,
 pdftitle={A Contract-Oriented Review of StradFixed2 and a Proposed Third Revision},
 pdfauthor={Technical analysis prepared for Vladimir Reshetnikov}}
\pagestyle{fancy}\fancyhf{}\fancyhead[L]{\small\textsc{Radical denesting: StradFixed2 review}}
\fancyhead[R]{\small\texttt{9000d6f}}\fancyfoot[C]{\thepage}
\setlength{\parskip}{4pt}\setlength{\parindent}{0pt}
\setlist{nosep,leftmargin=1.5em}
\lstset{basicstyle=\ttfamily\small,breaklines=true,columns=fullflexible,
 keepspaces=true,showstringspaces=false,frame=single,rulecolor=\color{black!15},
 backgroundcolor=\color{black!2},aboveskip=9pt,belowskip=9pt}
\newcommand{\Q}{\mathbb Q}\newcommand{\C}{\mathbb C}
\newcommand{\cost}{\operatorname{cost}}\newcommand{\depth}{\operatorname{depth}}
\newcommand{\Gal}{\operatorname{Gal}}\newcommand{\Nrm}{\operatorname{N}}
\newcommand{\code}[1]{\texttt{\detokenize{#1}}}
\newcommand{\src}[1]{\cite{source}, lines~#1}
\newtheorem{proposition}{Proposition}[section]
\newtheorem{theorem}[proposition]{Theorem}
\newtheorem{lemma}[proposition]{Lemma}
\theoremstyle{definition}\newtheorem{definition}[proposition]{Definition}
\theoremstyle{remark}\newtheorem{remark}[proposition]{Remark}
\newcolumntype{P}[1]{>{\raggedright\arraybackslash}p{#1}}
\begin{document}
\begin{titlepage}
\vspace*{15mm}
{\Large\color{ink}\textsc{Radical Denest / Technical Review}\par}
\vspace{11mm}
{\huge\bfseries A Contract-Oriented Review\par of \texttt{StradFixed2.wl}\par}
\vspace{6mm}
{\Large Correctness boundaries, search-state defects,\par resource contracts, and a proposed third revision\par}
\vspace{13mm}
{\large Prepared for Vladimir Reshetnikov\par}
{\large September 2026\par}
\vspace{12mm}
\begin{tabular}{@{}ll@{}}
Reviewed commit & \texttt{9000d6f533a68ee7659a6830ad25bba401f800a5}\\[3pt]
Reviewed file & \texttt{src/corrected/StradFixed2.wl}\\[3pt]
Baseline Git blob & \texttt{0bdeded347385a55a044fdd35883006d053685dc}\\[3pt]
Proposed package & \texttt{StradFixed3.wl}, context \texttt{RadicalDenest3`}\\[3pt]
Baseline integrity & 41,416 bytes; Git-blob hash verified\\[3pt]
Executed checks & 396 independent exact-mathematical and structural checks\\[3pt]
Native regression suite & 40 Wolfram tests supplied; not executed here
\end{tabular}
\vfill
\begin{minipage}{\textwidth}
\textbf{Reading the conclusions.} The strongest findings concern domain admission,
loss of previously accepted improvements, a missed index-reduction opportunity,
budget-sensitive negative caching, and resource-contract gaps. This review does
\emph{not} demonstrate an incorrect returned numerical value on the intended
closed exact-algebraic input class. The proposed package is a reproducible source
revision, not a claim of completed native-kernel qualification.
\end{minipage}
\end{titlepage}

\begin{abstract}
The second corrected Strad implementation has a valuable architecture: untrusted
candidate generation is separated from an exact-equality acceptance gate, and
symbolic hosts are processed through restricted congruence-preserving traversal.
The important remaining questions are therefore more specific than whether a
particular radical identity happens to simplify. Does the domain validator really
establish algebraicity? Can a later unsuccessful search discard a previously
certified improvement? Do recursive caches remember a mathematical result or
merely an exhausted budget? Are the advertised operation and batch limits applied
where the work is performed?

This article audits those questions against the pinned source, the unified-B
review, the repository survey, selected primary literature, and current Wolfram
documentation. It develops explicit witnesses for the sparse multiquadratic
ansatz and for index reduction that is useful without recursive progress. It
also proves a small-field character-reconstruction method and derives a general
odd-index trace--norm kernel. The accompanying package preserves the existing
quadratic, cubic, Honsbeek, Kummer, multiplier, rationalization, and traversal
machinery while hardening their interfaces. An exact copy of the baseline, a
hash-checking build script, a source diff, independent executable checks, and a
native regression suite make the proposal auditable and reproducible.
\end{abstract}
\begingroup
\small
\setlength{\parskip}{0pt}
\tableofcontents
\endgroup
\newpage

\section{Scope, evidence, and overall assessment}

\subsection{What was reviewed}
The source under review is the 745-line file \code{src/corrected/StradFixed2.wl}
at the commit printed on the title page. The copy included with this article
has the exact Git-blob SHA-1 reported by GitHub. The hash is computed using the
Git object convention, namely SHA-1 of \code{blob <byte-count>\0} followed by the
file bytes; it is not merely SHA-1 of the file content. The review is therefore
about a reproducible version, not a moving \code{main} branch. The preceding
review is \nolinkurl{src/code-review/unified-B/unified_analysis_B.tex}; the broader
survey is \nolinkurl{docs/report/radical_denesting_unified.tex}.\cite{source,prior,survey}

The literature reading concentrated on the repository's corrected editions of
Jeffrey--Rich, Borodin--Fagin--Hopcroft--Tompa, and Landau's 1989 extended
abstract. These sources are especially relevant to restricted practical
algorithms, square-root/fourth-root distinctions, partial denesting, and the
importance of specifying the base field and allowed roots of unity. The
repository explicitly identifies these files as editorial re-typesettings,
not facsimiles; this review does not silently attribute the editors'
qualifications to the original authors.\cite{jeffrey,bfht,landau,editions}

\subsection{Evidence levels}
The findings use four distinct kinds of evidence. \emph{Static evidence} is an
implication of the actual source: for example, an assignment overwrites an
incumbent regardless of its quality. \emph{Mathematical evidence} is an exact
identity or field-theoretic obstruction proved here. \emph{Documented semantic
evidence} comes from Wolfram's description of a function. \emph{Native execution}
would mean running this precise revision in a Wolfram kernel.

The first three kinds are available. The fourth is not: the Wolfram connector
was attempted but failed, and no local native kernel was available. Consequently,
none of the 40 supplied Wolfram tests is reported as passed. The 396 executed
checks are Python/SymPy exact mathematics, finite loop models, baseline-integrity
checks, and lexical/source checks. They are not a Wolfram emulator and do not
establish Wolfram evaluation semantics or performance.

This separation matters particularly for two results. The root-domain problem
is a defect in the predicate's implication, not an observed wrong numerical
answer. The lost-incumbent problem has a concrete control-flow witness and a
mathematically natural partial-denesting family, but the default package was
not run to measure which instances discover that partial denesting first.

\subsection{What the second correction gets right}
The dominant positive feature is the single acceptance gate. A proposed
replacement must be an explicit radical expression, fit the size limit, be
strictly cheaper than the current incumbent, and be certified equal to the
\emph{original target} of the island. This applies to specialized kernels,
\code{Simplify}, \code{Factor}, \code{ToRadicals}, and user callbacks. Searching
heuristically is entirely compatible with this architecture.\src{224--255}

The implementation also keeps \code{Root} objects from masquerading as denested
radical output, isolates \code{$Assumptions}, restricts traversal through
symbolic hosts, bounds the overall multiplier queue at its admission point,
and preserves only completed top-level passes on interruption. The polynomial
inverse used for denominator rationalization is algebraically correct. The
indirect square-root criterion and the principal phase for negative cube roots
are also correct; neither should be undone in another cleanup.\src{261--279,
301--323,379--396,648--690}

The appropriate conclusion is not that the package should be discarded.
It is that the separation between discovery and verification should be extended
to \emph{domain admission, incumbent state, cache state, and resource accounting}.

\section{Architecture and intended contracts}

\subsection{The execution pipeline}
\begin{longtable}{@{}P{.21\textwidth}P{.72\textwidth}@{}}
\toprule
Layer & Role in the baseline\\\midrule\endhead
Public entry points & \code{Strad}, \code{DenestRadicals}, \code{DenestCore}, and
\code{DenestReport} resolve per-head options and enter a shared session.\\
Session & A dynamically scoped configuration, deadline, statistics, trace,
certificate sample, memo table, and recursion depth are shared across the call.\\
Traversal & \code{walk} visits supported arithmetic/list hosts;
\code{island} handles exact algebraic subexpressions and recombination.\\
Discovery & Shape-specific square/cube kernels, linear factors in radical
extensions, index reduction, generic CAS proposals, and multiplier search
produce candidate presentations.\\
Acceptance & \code{accept} checks grammar, size, cost, and exact equality.
\code{acceptWithPolish} proposes several alternative presentations.\\
Publication & A completed pass is committed only after a whole-number equality
check when the entire input is numeric. Symbolic hosts rely on certified local
replacements and congruence.\\\bottomrule
\end{longtable}
This organization is visible directly in the source and is a substantial
improvement over a simplifier that assumes its candidate formulas always
choose the right branch.\cite{source,prior}

\subsection{Four separate invariants}
For an original island $E$ and incumbent $B$, the desired invariants are:
\begin{align}
& B=E, \label{eq:semantic-invariant}\\
& \cost(B)\leq_{\rm lex}\cost(E),\label{eq:cost-invariant}\\
& \text{every stage returns either }B\text{ or a certified cheaper candidate},
\label{eq:stage-invariant}\\
& \text{an exhausted search is not represented as a proof of nonexistence}.
\label{eq:search-invariant}
\end{align}
The baseline gate is designed to enforce the first two. The later-stage reset
violates the third while preserving equality. Budget-insensitive negative
memoization violates the fourth. Conflating these failures with incorrect
arithmetic would obscure both the diagnosis and the fix.

The cost is a presentation heuristic, not an algebraic invariant:
\begin{equation}
\cost(E)=\bigl(o(E),d(E),r(E),\operatorname{LeafCount}(E),b(E)\bigr).
\end{equation}
Here $o$ counts opaque algebraic representations, $d$ is radical depth, $r$ is
radical-node count, and $b$ prices exact coefficients. Equivalent algebraic
numbers can have very different costs. There is no contradiction in using
canonical algebraic values for equality and noncanonical radical expressions
for optimization.

\section{Finding register}
\small
\begin{longtable}{@{}P{.06\textwidth}P{.16\textwidth}P{.47\textwidth}P{.22\textwidth}@{}}
\toprule
ID & Priority & Finding and evidence & Revision status\\\midrule\endhead
F1 & High: contract & Exact-looking \code{Root} objects are admitted without
establishing algebraicity. Static and documented semantic evidence. & Conservative
rational-polynomial root grammar.\\
F2 & High: search state & Negative-radicand and Kummer stages can reset a better
incumbent to the original target. Static witness; partial-denesting example. &
Incumbent passed through both stages.\\
F3 & Medium: coverage & Index reduction is discarded unless recursion changes
its subproblem and reduces depth. Exact fourth-root witness. & Offer reduced
presentation before recursion.\\
F4 & Medium: cache & Unchanged results are cached without recursion/trial-budget
context. Static evidence. & Positive incumbents only; path-local cycle guard.\\
F5 & Medium: resources & Full integer factorizations and some expansion/
normalization work bypass the operation wrapper; scans and final costs lie
outside the outer region. & Explicit wrappers; budgeted scans and cached costs.\\
F6 & Medium: resources & The nominal 24-proposal discriminant batch can reach
46, although the global admission cap holds. Exact finite loop model. & Check
batch limit in the inner loop; configurable cap.\\
F7 & Medium: API semantics & A numerical diagnostic can produce the public
status \code{"Different"} despite exact-decision wording. Previously acknowledged
limitation, not a new discovery. & Numerical evidence is diagnostic only.\\
F8 & Low/medium: API & Malformed non-rule option tails can remain unevaluated
instead of reaching the documented \code{Failure} path. & Broader tail patterns,
central validation retained.\\
F9 & Low: metric & Gaussian-rational bit payloads are underpriced; subtractive
opaque-node counting is brittle. & Direct recursive accounting and opaque pruning.\\
F10 & Coverage opportunity & The sparse multisurd basis omits valid square-class
cosets; an explicit example proves the kernel's miss. Not a proof of whole-package
failure. & Small-rank character reconstruction added.\\\bottomrule
\end{longtable}
\normalsize

\section{The exact-algebraic boundary}

\subsection{Exact is not the same as algebraic}
The problematic clause is short:
\begin{lstlisting}
opaqueQ[e], TrueQ[Precision[e] === Infinity] && FreeQ[e, _Real]
\end{lstlisting}
It occurs inside \code{algebraicFormQ}, where \code{opaqueQ} includes both
\code{Root} and \code{AlgebraicNumber}. No polynomial or coefficient-domain check
is performed.\src{143--166}

Wolfram's \code{Root} is more general than a canonical algebraic-number wrapper.
Its documented representations include roots of general equations, polynomial
roots with nonrational numeric coefficients, and parameterized polynomial
roots. Exactness of representation does not establish membership in
$\overline{\Q}$.\cite{wlroot}

For example, consider an indexed representation of a root $r$ of
\begin{equation}\label{eq:trans-root}
 r^5+\pi r+1=0.
\end{equation}
The polynomial is strictly increasing on the real line, so it has a unique real
root, and that root is not zero. If $r$ were algebraic, then
\begin{equation}
 \pi=-\frac{r^5+1}{r}
\end{equation}
would be algebraic, a contradiction. Thus this is an exact number which is not
algebraic. An exact-looking indexed \code{Root[#^5 + Pi # + 1 &, 1]} retained by
the kernel satisfies the old structural acceptance rule without satisfying its
claimed mathematical conclusion. The native regression suite checks rejection
of this input; the present review does not claim to have observed its exact
internal form on a running kernel.

A parameterized \code{Root[#^5 + a # + 1 &, 1]} exposes a second issue: the input
need not even denote a closed numeric constant. The existing precision/no-
\code{Real} test cannot serve as a proof that parameters are absent.

\subsection{Why this matters even without a demonstrated wrong answer}
The exact method in \code{PossibleZeroQ} has a documented guarantee for
\emph{explicit algebraic numbers}. The source chooses that method after
\code{exactQ} succeeds. Therefore the validator is part of the proof boundary,
not a cosmetic input filter. A validator that admits a larger domain weakens
the justification for the subsequent certification call.\cite{wlzero}

The failure mode should nevertheless be stated narrowly. Unsupported roots may
simply remain unchanged, cause a timeout, or make the equality oracle return
\code{"Unknown"}. Those outcomes are plausible and safe. This review has not
shown that the baseline accepts a false equality on such an input. The finding
is that its promised implication
\begin{equation}
\code{ExactAlgebraicQ[e]}=\code{True}\ \Longrightarrow\ e\in\overline{\Q}
\end{equation}
is not established by the implementation.

\subsection{A conservative repair}
The revision accepts indexed univariate roots only when the defining function
produces a nonconstant polynomial with rational coefficients and the index is
an integer in the degree range. The supported optional root-isolation flag is
also checked. An \code{AlgebraicNumber} must have an admitted generator and a
Gaussian-rational coefficient list. Plus, Times, rational Power, and Gaussian
rationals retain their previous recursive treatment.

This is deliberately a \emph{sufficient grammar}. It rejects some genuinely
algebraic objects, including more general triangular representations and roots
with algebraic coefficients that have not been reduced to a rational-polynomial
form. Such inputs may first be normalized with a bounded \code{RootReduce} and
then rechecked. The revision does not claim to decide algebraicity of arbitrary
Wolfram expressions. Root-defining functions are ordinary Wolfram code, so this
validator is not a sandbox for hostile evaluation or side effects.

\subsection{Make all three equality statuses honest}
The baseline evaluates a difference numerically at 40 digits and, when its
reported accuracy and magnitude exceed thresholds, returns
\code{"Different"}. This branch never accepts a candidate, which is an important
safety distinction. But it is not the same as an exact proof of inequality, and
its existence conflicts with the public exact-decision wording.\src{186--212}

The proposed revision keeps the option for compatibility but changes its role.
Numerical evidence may increment a diagnostic statistic; it neither returns
\code{"Different"} nor prevents exact certification. Exact zero from
\code{RootReduce} proves equality. A rational or Gaussian-rational nonzero
normal form proves inequality. A validated polynomial root with nonzero
constant coefficient cannot be zero. Remaining cases go through the exact
algebraic method; interruption or uncertainty produces \code{"Unknown"}.
\cite{wlreduce,wlzero}

The record for an accepted candidate now identifies which exact equality route
succeeded. This is useful provenance, but it is still not an independent proof
object: both discovery and certification depend on the same kernel.

\begin{proposition}[Conditional soundness of the acceptance gate]
Assume that domain admission implies a finite exact algebraic value and that
all equality-oracle answers \code{"Equal"} are correct. If a gate returns the
incumbent unchanged or replaces it only by a cheaper radical expression
certified equal to the original target, then every sequence of successful
stage transitions preserves equality to that target and never increases cost.
\end{proposition}
\begin{proof}
The initial incumbent equals the target. An unchanged transition preserves both
properties. A replacing transition establishes equality to the original target
and strict cost decrease by its explicit tests. Induction proves the claim.
The hypotheses explain why domain validation and exact-oracle semantics belong
to the trusted core, whereas discovery formulas need not be trusted.
\end{proof}

\section{Search-state defects and concrete witnesses}

\subsection{Later no-op stages can erase earlier success}
In \code{powerProposals}, the assignments
\begin{lstlisting}
best = negativeRadicandCandidate[target, rho, p, q];
...
best = kummerCandidates[target, rho, p, q];
\end{lstlisting}
do not pass \code{best}. Both callees initialize or fall back to
\code{target}. A later no-op is thus interpreted as ``replace the incumbent by
the original input,'' rather than ``retain the incumbent.''\src{435--464,598--620}

An important qualification is that the specialized fast kernels usually return
depth-one candidates, after which \code{powerProposals} exits early. It would
therefore overstate the problem to claim that every successful quadratic
simplification is subsequently lost. The exposed case is a \emph{partial}
improvement, notably one found while separating a negative radicand.

Consider the positive numbers
\begin{align}
 t&=11-2\sqrt{29}, & w&=\sqrt{55-10\sqrt{29}},\\
 \rho&=16-2\sqrt{29}+2w, & B&=\sqrt5+\sqrt t.
\end{align}
Since $11>2\sqrt{29}$, $t>0$, and $w=\sqrt{5t}$. Hence
\begin{equation}\label{eq:partial}
 B^2=5+t+2\sqrt{5t}=\rho,\qquad B>0.
\end{equation}
This is the partial-denesting example discussed in the square-root
literature.\cite{bfht} The negative-radicand target is
\begin{equation}
 E=\sqrt{-\rho}=iB.
\end{equation}
Its original presentation has depth three, while $iB$ has depth two.
If recursive processing of the positive modulus discovers $B$, the negative
stage can accept $iB$. A subsequent Kummer stage which returns its original
target discards that certified improvement.

There is also a field-theoretic reason that a linear-factor search need not
rediscover the improvement. Put $F=\Q(\sqrt{29})$ and $K=F(w)$. This is a real
field, so $iB\notin K$. Moreover $B=\sqrt5(1+w/5)$, and $\sqrt5\notin K$:
$5$ is not a square in $F$, and $5/(5t)=1/t$ is not a square in $F$, since
$\Nrm_{F/\Q}(t)=5$ is not a rational square. Thus even the positive square root
is outside the displayed radicand field. This does not prove how every kernel
factorization is represented; it explains why retaining independent discoveries
is mathematically necessary.

The supplied white-box regression injects the partial modulus result and an
empty linear-factor result. It tests the stage interface directly without
pretending that default CAS discovery was observed. The repair is simply to
pass an incumbent through both functions and to return it on every no-op path.

\subsection{Useful index reduction does not require recursive progress}
The baseline takes a root $\gamma$ of $x^k-\rho$, writes $q=k\,r$, and recursively
processes $\gamma^{1/r}$. It offers the result only if recursion changed the
expression and reduced radical depth.\src{443--452}

Both requirements are too strong. Set
\begin{equation}\label{eq:index-witness}
 \rho=3+2\sqrt2=(1+\sqrt2)^2,\qquad
 E=\rho^{1/4}=\sqrt{1+\sqrt2}.
\end{equation}
The reduced expression already has a smaller presentation even though both
forms have radical depth two and two nonintegral-power nodes. The coefficient
and multiplication in $2\sqrt2$ disappear. With recursion disabled, the old
condition necessarily rejects this opportunity.

The example fits the Kummer search exactly. For $k=2$, the radical field
$\Q(\sqrt2)$ contains $\gamma=1+\sqrt2$. It does not contain the positive square
root of $\gamma$: the conjugate $1-\sqrt2$ is negative, while the conjugate of a
square in this real quadratic field must be nonnegative. Thus asking for a
linear factor of $x^4-\rho$ first does not make the $k=2$ step redundant.

The revised stage first offers the unmodified reduced presentation, with the
necessary root-of-unity phases and the original numerator exponent. It then
optionally recurses and offers those results too. The global acceptance gate,
not a separate depth-only predicate, decides whether each proposal is useful.
This also admits equal-depth improvements in node count or coefficient size.

\subsection{A failed bounded search is not a reusable negative fact}
The baseline memo key is the target expression alone. On completion it stores
\code{best}, including the case \code{best === target}; a later visit returns
that entry immediately. Recursive calls, however, use fewer remaining recursion
levels and divide the trial allowance by four. Local operation failures can
also occur without exhausting the session deadline.\src{430--433,628--639}

Suppose a target is first visited near the recursion limit. It may return
unchanged because the remaining search is weak, and that unchanged result is
cached. A later shallower visit has more opportunities, but the memo lookup
prevents them from being tried. The cache has confused ``this search found
nothing'' with ``nothing useful remains to be found.''

One repair is a resource-aware search-state key. Another is to store only
certified successes and use them as incumbents rather than terminal answers.
The revision chooses the latter. It also maintains a path-local in-progress
set, so removing negative entries does not invite immediate recursive cycles.
A revisit on the current path returns the best already known representation;
a later visit off that path is allowed to search again.

This policy can perform repeated unsuccessful work. That is a conscious tradeoff
for simpler correctness of search-state reuse. A more advanced implementation
could store an incumbent plus a resumable frontier and a clearly specified
coverage descriptor, but an undifferentiated memoized expression is insufficient.

\section{Resource contracts, option dispatch, and metric details}

\subsection{An overall timeout does not enforce every local timeout}
The header promises a bounded region for every expensive kernel operation.
The source nevertheless performs \code{FactorInteger} directly while building
multisurd bases and discriminant-prime lists. It also expands an ansatz square
and normalizes some extracted linear roots outside \code{bounded}.\src{357,
364,425--426,564}

These calls are usually inside the \emph{outer} session constraint. Therefore
it would be incorrect to say that the code has no time limit at all. The defect
is that \code{"OperationTime"} does not govern the operations in question.
A full factorization of a large discriminant can consume most of the session,
even though only eight resulting primes will be retained. Taking a prefix
\emph{after} complete factorization does not bound factorization effort.

The revision wraps these explicit calls and treats failure as absence of that
proposal source. It also bounds the combined shape-specific proposal call and
rejects oversized candidates before polishing. These are practical safeguards,
not a proof that expression evaluation has constant overhead or that every
internal arithmetic step is separately interruptible.

\subsection{The discriminant batch bound is misplaced}
In the old prime-power loop, the batch counter is checked before entering each
prime, but not before each exponent:
\begin{lstlisting}
Do[If[batch >= 24, Break[]];
   Do[batch++; admit[m prime^e], {e, 1, q - 1}],
   {prime, primes}]
\end{lstlisting}
For $q=24$, one prime contributes 23 proposals. Because $23<24$, a second prime
is entered and contributes another 23. The batch reaches 46. For $q=32$, the
first prime alone reaches 31. Among $2\le q\le32$ and up to eight primes, the
largest such overshoot is 46. The independent loop-model checks exhaust this
finite range.\src{565--572}

The queue's separate admission guard still prevents more than
\code{"MultiplierCap"} entries from being admitted. It is the local generation
budget, not the global queue bound, that fails. The revision checks all relevant
bounds in the exponent loop and exposes the old constant as
\code{"DiscriminantBatchCap"}, defaulting to 24. A value of zero disables that
expansion source without disabling seed multipliers.

\subsection{Budgeted work should include preflight and report metrics}
The original outer region contains the pass loop, but input scans before it and
\code{RadicalCost} evaluations used to assemble the final report lie outside
it. A zero time budget therefore still permits full cost traversal of the input.
For large trees, this contradicts a literal whole-call budget.\src{692--728}

In the revision, input scans and initial/candidate costs are performed inside
the outer region. The committed snapshot carries its already-computed cost.
When the initial cost was never computed, the report uses
\code{Missing["NotComputed"]} instead of performing expensive post-timeout work.
The zero-budget path does not scan the expression merely to populate cosmetics.

Ordinary argument evaluation, option resolution, final association construction,
and front-end serialization are not claimed to be covered. A real hard boundary
requires a parent process supervising a worker kernel; the parent can enforce
wall-clock and operating-system memory limits and recover a previously serialized
checkpoint. That is an architectural extension, not something accomplished by
adding another \code{TimeConstrained} around the same kernel expression.

Wolfram documents interrupt granularity and abort protection for time constraints,
and states that their accounting concerns main-kernel CPU time rather than all
threads/processes. Memory constraints are likewise not a process-RSS sandbox.
The revision's report therefore describes its limits as cooperative kernel
limits.\cite{wltime,wlmemory}

\subsection{Degree limits are output filters, not complexity certificates}
The package computes a minimal polynomial and then checks its degree against
\code{"MaxDegree"}. That is a sensible acceptance filter, but a polynomial of
unacceptable degree may already have been expensive to compute. The same
observation applies to extension construction and algebraic normalization.
The revision retains this behavior rather than falsely promising a cheap
pre-computation degree bound.

A future preflight may multiply upper bounds for independent radical degrees,
reduce known square-class relations, and stop when the bound becomes too large.
Such a bound must be conservative: naive multiplication often dramatically
overestimates the true field degree. A pessimistic bound is useful for resource
policy, not a reason to claim that a denesting is mathematically unavailable.

\subsection{Malformed options should reach the validator}
The validator checks that option items are rules and returns a structured
\code{Failure}. But the baseline public definitions use \code{OptionsPattern[]}
for the tail. A non-rule argument such as 123 need not match those definitions,
so it may remain an unevaluated public call rather than reaching the validator.
This is a dispatch gap between the documented API and its error path.
\src{120--138,736--742}

The proposed public definitions accept an arbitrary option tail and pass it to
the same validator. The more specific legacy Boolean overload remains. Unknown
rule names, invalid values, and malformed non-rule tails consequently share a
consistent error mechanism. Missing required positional arguments remain an
ordinary arity issue; this change does not purport to define every malformed
call imaginable.

\subsection{Cost accounting should be direct}
The original radical-node counter counts all powers and then subtracts powers
inside opaque algebraic objects. A pruned recursive traversal is simpler and
cannot double-subtract a radical if an opaque representation contains another
opaque object. This review treats the latter as a representation-dependent
robustness issue, not a demonstrated negative count for a particular native
canonical input.\src{173--181}

There is a more direct coefficient issue: exact complex numbers are atomic to
ordinary expression traversal. Searching for integer/rational atoms therefore
does not reliably price their real and imaginary components. The replacement
\code{bitSize} handles \code{Complex} explicitly, then recurses through other
non-atomic payloads. Opaque counts and radical-node counts prune opaque bodies;
bit size still prices their stored data. The revised cost documentation states
this distinction rather than leaving it implicit.

Even a perfectly implemented lexicographic metric is noncompositional under
all Wolfram evaluations. A local simplification can change sharing, cancellation,
or a parent's final representation. The package still checks whole-pass cost
before commitment. A future DAG-aware metric or a small Pareto frontier could
improve coverage, but neither is implemented here.

\section{Verification of the retained algebraic kernels}

This section records the mathematical contracts that the implementation should
preserve. It also separates a formula being a valid \emph{proposal} from it
already being the correct principal value. The exact gate remains the final
authority for the latter.

\subsection{Direct square-root denesting}
Let $a,b,c\in\Q$, $c>0$, $\sqrt c\notin\Q$, and suppose
$\rho=a+b\sqrt c>0$. Put $D=a^2-b^2c$. When $D\ge0$ and $s=\sqrt D\in\Q$,
set
\begin{equation}
 u=\frac{a+s}{2},\qquad v=\frac{a-s}{2}.
\end{equation}
Under $a>0$, both $u$ and $v$ are nonnegative, because $s\le a$. Moreover
$uv=b^2c/4$, so
\begin{equation}
 \left(\sqrt u+\operatorname{sgn}(b)\sqrt v\right)^2
 =a+b\sqrt c.
\end{equation}
Since $u\ge v$, the displayed candidate is nonnegative even when $b<0$.
It is therefore the principal square root. The source absorbs $|b|$ into the
radicand before applying the recipe, which changes its internal representation
but not this argument. The norm/discriminant mechanism is the familiar
quadratic denesting criterion discussed in the literature.\cite{bfht,jeffrey}

\subsection{Indirect denesting with a fourth root}
Suppose $D<0$ and $b>0$. Define
\begin{equation}
 h=\sqrt{b^2-a^2/c}.
\end{equation}
When $h\in\Q$, let $u=(b+h)/2$ and $v=(b-h)/2$. Here $0\le h\le b$, so
$u,v\ge0$, and $uv=a^2/(4c)$. Consequently,
\begin{align}
 \left[c^{1/4}\left(\sqrt u+\operatorname{sgn}(a)\sqrt v\right)\right]^2
 &=\sqrt c\left(b+\frac a{\sqrt c}\right)\\
 &=a+b\sqrt c.
\end{align}
The candidate is nonnegative because $u\ge v$. This proves the source's recipe.
The rational-square condition is
\begin{equation}
 h^2=-D/c,\qquad
 h\in\Q\ \Longleftrightarrow\ -cD\text{ is a rational square}.
\end{equation}
It is not the condition that $-D$ alone be a rational square. For instance,
\begin{equation}
 \sqrt{4+3\sqrt2}=2^{1/4}(1+\sqrt2),
\end{equation}
where $D=-2$ and $-cD=4$. This distinction is already emphasized in the survey
and preceding correction and is retained here.\cite{survey,prior,bfht}

\subsection{Principal roots on the negative real axis}
Wolfram rational powers use a principal complex value. For $r>0$ and rational
$p/q$ in lowest terms,
\begin{equation}
 (-r)^{p/q}=e^{i\pi p/q}r^{p/q}=(-1)^{p/q}r^{p/q}.
\end{equation}
For odd $q$, this generally differs from the signed real $q$th root convention.
In particular,
\begin{equation}
 (-7-5\sqrt2)^{1/3}=(-1)^{1/3}(1+\sqrt2),
\end{equation}
not $-1-\sqrt2$. A power identity alone does not choose the correct root:
all $q$ root-of-unity multiples have the same $q$th power. The implementation's
phase enumeration and final exact comparison against the original principal
expression are therefore essential.

Similarly, unrestricted transformations such as
$(uv)^{1/q}=u^{1/q}v^{1/q}$ or $(u^k)^{1/k}=u$ must not be made into trusted
rewrite rules. They can be used to organize a finite candidate orbit whose
members are checked individually.

\subsection{Odd roots in a quadratic field}
The cubic kernel admits a clean generalization. Suppose
\begin{equation}
 \beta=A+B\sqrt c,\qquad A,B,c\in\Q,
\end{equation}
with conjugate $\bar\beta=A-B\sqrt c$. Define the trace and norm
\begin{equation}
 T=\beta+\bar\beta=2A,\qquad n=\beta\bar\beta=A^2-cB^2.
\end{equation}
Both roots satisfy $x^2-Tx+n=0$. Thus the polynomial sequences
\begin{align}
 D_0&=2,&D_1&=T,&D_j&=T D_{j-1}-nD_{j-2},\\
 S_0&=0,&S_1&=1,&S_j&=T S_{j-1}-nS_{j-2}
\end{align}
satisfy
\begin{equation}\label{eq:D-S-meaning}
 D_j=\beta^j+\bar\beta^j,\qquad
 (\beta-\bar\beta)S_j=\beta^j-\bar\beta^j.
\end{equation}
These identities follow by the same second-order recurrence and the initial
values. They also give the polynomial identity
\begin{equation}\label{eq:D-S-identity}
 D_j(T,n)^2-(T^2-4n)S_j(T,n)^2=4n^j.
\end{equation}
It is a polynomial identity even when the two quadratic roots coincide, since
it holds for indeterminate distinct roots and hence identically.

Now let $q\ge3$ be odd and suppose we seek a root inside $\Q(\sqrt c)$ of
\begin{equation}
 \beta^q=a+b\sqrt c,\qquad b\ne0.
\end{equation}
Taking norms gives the necessary condition
\begin{equation}
 n^q=a^2-b^2c.
\end{equation}
The rational candidate for $n$ is the signed real $q$th root of the right-hand
side. For every rational root $T$ of
\begin{equation}\label{eq:trace-polynomial}
 D_q(T,n)-2a=0,
\end{equation}
form
\begin{equation}\label{eq:odd-reconstruction}
 \beta=\frac T2+\frac b{S_q(T,n)}\sqrt c.
\end{equation}

\begin{proposition}[Trace--norm reconstruction]
Under the assumptions above, if $n,T\in\Q$ satisfy the norm equation and
\eqref{eq:trace-polynomial}, then $S_q(T,n)\ne0$ and the value in
\eqref{eq:odd-reconstruction} satisfies $\beta^q=a+b\sqrt c$.
Conversely, every such root in $\Q(\sqrt c)$ supplies such a pair $(T,n)$.
\end{proposition}
\begin{proof}
Substituting the equations for $D_q$ and $n^q$ into
\eqref{eq:D-S-identity} yields
\begin{equation}
 (T^2-4n)S_q(T,n)^2=4b^2c\ne0.
\end{equation}
In particular $S_q\ne0$. Writing $A=T/2$ and $B=b/S_q$ gives
$A^2-cB^2=n$. Therefore the reconstructed quadratic pair has trace $T$ and
norm $n$, and \eqref{eq:D-S-meaning} shows
$\beta^q=D_q/2+B\sqrt c\,S_q=a+b\sqrt c$.
The converse follows by taking the trace and norm of the assumed root.
\end{proof}

For $q=3$,
\begin{equation}
 D_3=T^3-3nT,\qquad S_3=T^2-n,
\end{equation}
which recovers the existing cubic implementation. For $q=5$,
\begin{equation}
 D_5=T^5-5nT^3+5n^2T,\qquad
 S_5=T^4-3nT^2+n^2.
\end{equation}
The revision uses rational linear factors of the trace polynomial rather than
trying to solve an arbitrary quintic or septic by radicals. By default the
extension is enabled through index seven. This is a bounded additional proposal
kernel, not a completeness statement for all odd radicals: roots requiring a
larger field are outside its scope. The original cubic kernel is retained.

\subsection{The Honsbeek quartic proposal is algebraically sound}
Let $A^3=a$, $B^3=b$, and set
\begin{equation}
 N=B^2+sAB-\frac{s^2A^2}{2},\qquad D=b-s^3a.
\end{equation}
Expanding and reducing $A^3=a$, $B^3=b$ gives
\begin{align}
 N^2&=(b-s^3a)B+A\left(2sb+\frac{s^4a}{4}\right),\\
 N^2-(A+B)D
 &=\frac A4\left[a(s^4+4s^3)+b(8s-4)\right].\label{eq:honsbeek}
\end{align}
Thus every rational solution of
\begin{equation}
 s^4+4s^3+8(b/a)s-4b/a=0
\end{equation}
with $a\ne0$ and $D\ne0$ gives candidates $\pm N/\sqrt D$ whose squares equal
$A+B$. The source restricts to $D>0$ and offers both signs. The exact gate
selects the original principal square root.\src{325--343}

For $A=\sqrt[3]5$, $B=-\sqrt[3]4$, one may take $s=-2$, giving $D=36$ and
\begin{equation}
 \sqrt{\sqrt[3]5-\sqrt[3]4}
 =\frac{\sqrt[3]2+\sqrt[3]{20}-\sqrt[3]{25}}3.
\end{equation}
The positive sign is immediate because $\sqrt[3]2+\sqrt[3]{20}>3$ while
$\sqrt[3]{25}<3$. This familiar example is also present in the repository
literature.\cite{jeffrey,landau} The independent test suite verifies
\eqref{eq:honsbeek} symbolically in a polynomial quotient ring, rather than
merely checking decimal approximations.

\subsection{Polynomial inverses for rationalization}
If a nonzero algebraic denominator $d$ satisfies
\begin{equation}
 a_0+a_1d+\cdots+a_nd^n=0,\qquad a_0\ne0,
\end{equation}
then
\begin{equation}
 d^{-1}=-\frac{a_1+a_2d+\cdots+a_nd^{n-1}}{a_0}.
\end{equation}
This is exactly the Horner expression used in \code{rationalizeRaw}. The source
checks that coefficients are rational, excludes zero constant coefficient,
and verifies $d\cdot d^{-1}=1$ before using the inverse. There is no missing sign
or off-by-one error in the formula.\src{261--279}

The result need not be cheaper, which is why the standalone rationalization
API has a different cost contract from a denesting proposal. When a symbolic
numerator is present, multiplying by a certified inverse of a fixed nonzero
numeric denominator is a congruence step; it is not a license to cancel arbitrary
symbolic factors whose zero sets were not tracked.

\section{The sparse multisurd ansatz: a proved limitation}

\subsection{A counterexample inside a small field}
The baseline multisurd kernel forms two trial sets: 1 together with visible
prime factors, and then those values together with the visible radicands. It
solves for a rational linear combination of the corresponding square roots.
These are useful small ansatzes, but they do not constitute a basis for all
possible denestings.\src{345--375}

Take
\begin{equation}\label{eq:beta-example}
 B=\sqrt6+\sqrt{10}+\sqrt{21}.
\end{equation}
Direct expansion gives
\begin{equation}\label{eq:rho-example}
 B^2=37+4\sqrt{15}+6\sqrt{14}+2\sqrt{210}.
\end{equation}
The old largest ansatz uses
\begin{equation}
 U=\{1,2,3,5,7,14,15,210\}.
\end{equation}
It contains none of the square classes 6, 10, or 21. Distinct square classes
have linearly independent square roots over $\Q$ inside a multiquadratic field.
Consequently neither $B$ nor $-B$ lies in the rational span of
$\{\sqrt u:u\in U\}$. Since the only roots of $x^2-B^2$ are $\pm B$, the ansatz
has no solution.

For completeness, the independence statement follows from character
orthogonality. Choose independent square classes generating a field containing
all the roots. Each root belongs to a distinct one-dimensional eigenspace for
the sign-changing automorphism group. Applying the projector for one character
to a rational linear relation isolates its coefficient, which must therefore
vanish. This argument also motivates the replacement kernel below.

The conclusion is about the \emph{multisurd kernel}. Kummer, generic CAS, or
multiplier proposals might still solve the same whole input. This review does
not turn a demonstrated local search-space omission into an unmeasured
end-to-end failure claim.

\subsection{Closing the visible prime basis is still not enough}
A tempting repair is to include every product of the visible radicand primes.
That helps, but it does not solve the general issue. For example,
\begin{equation}
 \left(\sqrt{10}+\sqrt{15}+\sqrt{35}\right)^2
 =60+10\sqrt6+10\sqrt{14}+10\sqrt{21}.
\end{equation}
The prime 5 is absent from the radicands 6, 14, and 21. It occurs in coefficients,
and the square root lies in a square-class coset not contained in the field
obtained from those radicands alone.

In \eqref{eq:rho-example}, the radicand field is
$K=\Q(\sqrt{14},\sqrt{15})$. The classes supporting $B$ lie in
\begin{equation}
 6\langle14,15\rangle\subset\Q^\times/\Q^{\times2},
\end{equation}
not necessarily in the subgroup itself. Multipliers can sometimes discover
this missing coset; a direct reconstruction method can find it without first
guessing its representative.

\section{A bounded character-reconstruction kernel}

\subsection{Square-class basis without an explicit prime-factor basis}
Suppose
\begin{equation}
 \rho=a_0+\sum_j a_j\sqrt{d_j},\qquad a_j\in\Q,\quad d_j\in\mathbb Z_{>0}.
\end{equation}
Let $H$ be the subgroup of $\Q^\times/\Q^{\times2}$ generated by the $d_j$,
and let $r$ be its rank. Then $K=\Q(\sqrt h:h\in H)$ has degree $m=2^r$.

The new kernel constructs representatives of $H$ incrementally. Starting with
$\{1\}$, a new radicand $d$ is already represented precisely when $d/h$ is a
rational square for some existing representative $h$. Otherwise $d$ becomes a
new generator and the representative list is doubled by multiplication by $d$.
The rank cap is checked before that doubling.

There is no explicit \code{FactorInteger} call in this basis construction.
The Wolfram implementation uses exact rational-square tests involving square
roots; it does not promise that the kernel never factors internally while
simplifying them. An integer-square-root test on the numerator and denominator
would be a possible further low-level optimization.

\subsection{Conjugates and the finite transform}
Index $H$ and the automorphisms of $K$ by $\mathbb F_2^r$. Let
\begin{equation}
 \chi_g(\sigma)=(-1)^{g\cdot\sigma},\qquad
 \lambda_\sigma=\sqrt{\sigma(\rho)}
\end{equation}
with the principal square root at each embedding. Enumerate signs
$\epsilon_\sigma\in\{\pm1\}$, fixing $\epsilon_0=1$. For each choice compute
\begin{equation}\label{eq:fourier}
 F_g=\frac1m\sum_\sigma\chi_g(\sigma)\epsilon_\sigma\lambda_\sigma.
\end{equation}
If every $F_g^2$ is rational, each $F_g$ can be represented by one signed square
root of a rational number, with the sign determined by exact comparison. The
candidate is then
\begin{equation}\label{eq:reconstruct}
 \sum_g F_g.
\end{equation}
Orthogonality of the characters shows that this sum is exactly $\lambda_0$:
all columns except the identity column cancel. Nevertheless, the implementation
also runs the final exact equality gate. A proof of the algebraic construction
should not become a reason to bypass a robust implementation boundary.

\subsection{Why this captures the missing coset}
\begin{lemma}[Single-coset support]
Let $K/\Q$ be multiquadratic. Suppose $\beta\ne0$, $\beta^2\in K$, and $\beta$
lies in some finite multiquadratic extension of $\Q$. Then
$\beta\in\sqrt d\,K$ for some $d\in\Q^\times$.
\end{lemma}
\begin{proof}
Choose a finite multiquadratic field $M$ containing both $K$ and $\beta$.
For $\tau\in\Gal(M/K)$, the equality
$(\tau\beta)^2=\beta^2$ gives $\tau\beta=\pm\beta$. These signs form a character
of $\Gal(M/K)$, because $\beta\ne0$. The Galois group of $M/\Q$ is a finite
vector space over $\mathbb F_2$, so that character extends to the whole group.
Its one-dimensional eigenspace is generated by $\sqrt d$ for a rational square
class $d$. The quotient $\beta/\sqrt d$ is fixed by $\Gal(M/K)$ and hence lies
in $K$.
\end{proof}

\begin{theorem}[Restricted completeness of exhaustive reconstruction]
Ignore computational limits and suppose exact algebraic arithmetic and equality
are available. For $\rho\in K$, the sign enumeration in \eqref{eq:fourier}
finds a representation of the principal $\sqrt\rho$ as a sum of square roots of
rational numbers whenever such a representation exists. Signed rational
radicands, and thus complex square roots, may be allowed. This statement does
not include arbitrary fourth roots or higher-index radicals.
\end{theorem}
\begin{proof}
The zero case is trivial. A finite sum of rational square roots lies in a finite
multiquadratic extension. By the lemma, $\beta=\sqrt\rho=\sqrt d\,u$ with $u\in K$.
If $\sqrt d\in K$, absorb it into $u$ and use $d=1$. Otherwise each automorphism
of $K$ extends to $K(\sqrt d)$ while fixing $\sqrt d$.

Expand $u=\sum_g c_g\sqrt{h_g}$ in the square-class basis of $K$, where
$c_g\in\Q$. At each embedding, the image of $\beta$ is one of the two square
roots of $\sigma(\rho)$, so it equals
$\epsilon_\sigma\lambda_\sigma$ for some enumerated sign choice. At the identity
embedding the sign is $+1$. Character projection then gives
$F_g=c_g\sqrt{d h_g}$, and $F_g^2=c_g^2dh_g\in\Q$.
Thus that sign choice passes the rational-square tests and reconstructs $\beta$.
Conversely, whenever the tests pass, selecting the exact signs of the square
roots gives the $F_g$ themselves; orthogonality proves
\eqref{eq:reconstruct} equals the principal target.
\end{proof}

The theorem is a mathematical guarantee for an exhaustive restricted procedure,
not a guarantee for a timed execution. It is also not a new general denesting
theorem: it is an explicit character-based formulation suitable for this
implementation. The preceding review already pointed toward field bases and
character projections; this revision turns a small-rank form into executable
proposal code.\cite{prior}

\subsection{An explicit reconstruction table}
For \eqref{eq:beta-example}, write $a=\sqrt6$, $b=\sqrt{10}$, $c=\sqrt{21}$.
Choose generators 14 and 15 and extend their sign changes while fixing
$\sqrt6$. The principal square roots of the four conjugates are
\begin{equation}
 (\lambda_0,\lambda_1,\lambda_2,\lambda_3)
 =(a+b+c,\ a+b-c,\ a-b+c,\ -a+b+c).
\end{equation}
All four displayed values are positive. The only nontrivial triangle inequality
is $a+b>c$, which follows from $(a+b)^2=16+4\sqrt{15}>21$; the other two follow
already from $a^2+c^2>b^2$ and $b^2+c^2>a^2$.

Use the sign vector $(1,1,1,-1)$. The Hadamard transform gives
\begin{equation}
 (F_0,F_1,F_2,F_3)=(\sqrt6,\sqrt{21},\sqrt{10},0),\qquad
 (F_0^2,F_1^2,F_2^2,F_3^2)=(6,21,10,0).
\end{equation}
The missing square classes are recovered directly. With this generator order
and the implementation's binary sign enumeration, this is the second sign
pattern, not an expensive search through all possible rational multipliers.
The independent checks verify this example and 48 additional exact rank-two
coset examples using separate multiquadratic arithmetic and certified rational
intervals for signs.

\subsection{The cost is exponential in field degree}
There are $2^{m-1}=2^{2^r-1}$ sign patterns. The distinction between rank and
field degree is operationally important:
\begin{center}
\begin{tabular}{@{}rrr@{}}\toprule
Square-class rank $r$ & Field degree $m$ & Sign patterns\\\midrule
1 & 2 & 2\\2 & 4 & 8\\3 & 8 & 128\\4 & 16 & 32,768\\\bottomrule
\end{tabular}
\end{center}
The revision defaults to rank two and eight patterns. Rank three with 128
patterns is an explicit opt-in; rank four is accepted only with a deliberately
chosen pattern budget and remains potentially expensive. Each algebraic
reduction and sign certificate is also subject to the shared limits.

The older sparse ansatz remains as a fallback. Preserving it avoids throwing
away a useful fast heuristic for inputs outside the new rank cap. The new
kernel supplements rather than replaces the larger Kummer and multiplier
machinery.

\section{The proposed source revision}

\subsection{A standalone file, not an undocumented patch layer}
\code{StradFixed3.wl} is a complete loadable package in
\code{RadicalDenest3`}. It does not require loading version two first and does
not modify the user's repository. The archive also includes
\code{StradFixed2.original.wl}, the exact reviewed baseline,
\code{revision_definitions.wl}, and \code{build_revision.py}.

The builder verifies the baseline's Git-blob identity, applies anchored changes,
removes replaced function definitions, appends the new private definitions
before closing the package, and writes the standalone file. Every replacement
anchor has an expected occurrence count, so the builder fails rather than
silently applying a review to a different source version. A unified diff is
included for human inspection.

The original low-degree and general-search algorithms remain present. Additional
changes include fresh multisurd coefficient symbols, early size checking before
polishing, exact equality-method labels in accepted-proposal records, and
explicit report wording about cooperative resource limits and certificate scope.
These are small but useful measures against accidental coupling and misleading
diagnostics.

\subsection{Behavior intentionally changed}
The changes are not perfectly backward compatible. A valid algebraic value in
an unsupported \code{Root} representation can now be rejected by the conservative
grammar. Numeric prefiltering defaults to false, and setting it true no longer
produces an unproved inequality status. Gaussian bit accounting and opaque
pruning can change cost tie-breaks. A zero-budget report can contain missing
cost fields. Malformed option tails now produce structured failures instead
of remaining unevaluated.

These changes are documented because they are contract corrections, not
incidental implementation details. Existing callers that compare exact output
syntax should instead check equality and the desired presentation property,
unless a specific normal form is part of their own API.

\subsection{New options}
\begin{longtable}{@{}P{.31\textwidth}P{.12\textwidth}P{.50\textwidth}@{}}
\toprule
Option & Default & Meaning\\\midrule\endhead
\code{"MaxCharacterRank"} & 2 & Maximum independent square-class rank for the
new reconstruction kernel. Zero disables it; values above four are rejected.\\
\code{"MaxCharacterPatterns"} & 8 & Maximum sign patterns attempted. Zero
disables character reconstruction. A truncated enumeration is not a negative
mathematical answer.\\
\code{"MaxOddIndex"} & 7 & Largest odd index admitted by the additional
quadratic-field trace--norm kernel. The existing cubic path is retained.\\
\code{"DiscriminantBatchCap"} & 24 & Maximum prime-power/mixed-product proposal
attempts generated by one discriminant expansion. The global queue and trial
limits remain separate.\\\bottomrule
\end{longtable}
All existing options remain. \code{"MaxTrials"} continues to be per island,
not a global counter for the entire call; recursive invocations continue to use
smaller trial allowances. The overall session deadline and memory constraint
remain shared. Those distinctions are intentional and should not be reported
as regressions.

\subsection{Usage}
\begin{lstlisting}
Get["StradFixed3.wl"];

RadicalDenest3`DenestCore[
  Sqrt[37 + 4 Sqrt[15] + 6 Sqrt[14] + 2 Sqrt[210]],
  "MaxTrials" -> 0
]

RadicalDenest3`DenestReport[
  (41 + 29 Sqrt[2])^(1/5),
  "Trace" -> True,
  "MaxTrials" -> 0
]

RadicalDenest3`DenestRadicals[
  expression,
  "AllLevels" -> True,
  "MaxCharacterRank" -> 3,
  "MaxCharacterPatterns" -> 128
]
\end{lstlisting}
The first call targets $\sqrt6+\sqrt{10}+\sqrt{21}$ and the second targets
$1+\sqrt2$. These are mathematical expected results and native regression
fixtures; they are not a transcript of a kernel session performed for this
review. Fully qualified names are useful when loading versions two and three
side by side.

\section{Validation, reproducibility, and remaining qualification}

\subsection{What actually ran}
The independent suite completed 396 checks with zero failures under Python
3.13.5 and SymPy 1.14.0. The JSON output records every check. Its coverage includes:
\begin{itemize}
\item the baseline Git-object hash; symbolic Honsbeek, trace--norm, and polynomial
inverse identities; exact odd-root reconstruction instances;
\item exact square-class arithmetic, the sparse-ansatz witness, Hadamard
orthogonality, and exhaustive rank-two reconstruction witnesses;
\item the finite discriminant-batch model, separate global-admission caps,
partial-denesting/index-reduction identities, and lexical/source checks.
\end{itemize}
Polynomial identities use exact algebra. Sign isolation for the multiquadratic
examples uses rational outward bounds derived from integer square roots on a
$10^{-60}$ grid. No floating-point agreement is used as a proof of equality.
The word ``passed'' applies to these particular checks only.

The source checks verify balanced delimiters after removing strings and nested
comments and test for the intended critical source patterns. They cannot detect
all Wolfram semantic errors. In particular, they do not establish option
precedence, message behavior, native factorization choices, interrupt behavior,
or actual end-to-end runtime.

\subsection{The native regression suite}
The 40 supplied Wolfram tests cover conservative root admission, exact equality,
wrong branches, Gaussian coefficient cost, malformed options, disabled budgets,
quadratic/cubic/fifth-root examples, the missing-coset example, rationalization,
restricted host traversal, callback rejection, index reduction without recursive
progress, incumbent preservation, negative-cache removal, cycle cutting, and
certificate replay. Some are black-box tests; others deliberately stub discovery
functions so that a control-flow contract can be tested independently of CAS
heuristics.

Run them in a native kernel with
\begin{lstlisting}
wolframscript -file run_tests.wls
\end{lstlisting}
The runner writes \code{native_results.json} containing the kernel version and
success/failure counts, and exits with a nonzero status on failed tests. That
file is deliberately absent from this delivery: manufacturing it would imply
an execution that did not occur.

The tests are a starting qualification suite, not a claim that all behavior of
the preceding implementation has been covered. The repository's original
\code{StradFixed2Tests.wlt} should also be run against the revision with public
contexts adapted, treating intentional API changes separately. Prior reported
test success belongs to the prior implementation and environment; it cannot be
inherited by a source revision.\cite{prior}

\subsection{What should be measured next in a native environment}
The first native run should inspect all messages, not merely count equality
checks. A candidate that simplifies correctly only after emitting an unexpected
message may be turned into \code{$Failed} by a surrounding \code{Check}.
Deterministic tests should then be followed by generated exact families, varying
signs, numerator exponents, scalar factors, and repeated occurrences at different
recursion depths.

For timing, record wall time and the package's counters separately. Compare cold
and warm kernels because algebraic and factorization caches affect search cost.
Measure default rank-two character reconstruction before raising the rank cap.
Keep equality-preservation success distinct from denesting coverage: a result
can be correct yet unchanged, or correctly partially denested without reaching
the best known presentation.

The honest release status is therefore \emph{proposed research revision with
independently checked mathematics and supplied native tests}. It is not yet a
production-qualified replacement, and it is not a proof that the new defaults
outperform the old defaults on every input family.

\section{Further improvements, in priority order}

\subsection{Make the trusted core smaller and more explicit}
The domain recognizer, equality oracle, and commit mechanism should be separable
from discovery. A future certificate structure could contain the exact input
and output values, an annihilating polynomial, isolating data for the selected
embedding, and an independently replayable equality witness. A record saying
which kernel function returned true is useful audit metadata, but it is not
such a certificate.

A structural proof graph would distinguish accepted proposals, local replacements,
congruence rebuilding, and committed passes. The current bounded certificate
sample can contain proposals from an uncommitted pass. Its revised description
says so explicitly. It should not be interpreted as the exact derivation of the
published final result.

\subsection{Cache mathematical state separately from search coverage}
Positive-only caching avoids the immediate defect but leaves performance on the
table. A more capable cache entry would store a certified incumbent plus a
frontier of untried multipliers, divisors, and sign patterns. It could also record
which field presentation was used, which algorithms completed, and what resources
were available. An interrupted factorization is not a completed negative
field-membership test, and an exhausted prefix of sign patterns is not an
exhausted search space.

Canonical algebraic keys could improve deduplication, but normalizing every
cheap multiplier with \code{RootReduce} may dominate the search. One reasonable
compromise is cheap syntactic deduplication first and bounded algebraic
canonicalization only for promising expensive trials. No such optimization
should be justified by calling an expanded-expression string a canonical
algebraic value.

\subsection{Separate field discovery from radical presentation}
The current radical-extension list can contain redundant generators. A reduced
primitive-element representation, together with maps back to display radicals,
would make repeated field operations less wasteful. Square-class rank reduction
is the simplest instance; higher-index Kummer classes require a more general
representation.

A low minimal-polynomial degree is useful evidence that a presentation might
be obtainable, but it is not itself a denesting formula or a completeness
certificate. Similarly, \code{ToRadicals} is a proposal mechanism with documented
scope, not an unrestricted denesting oracle.\cite{wlradicals} The repository's
survey and Landau discussion emphasize why roots of unity and the chosen base
field cannot be treated as free without changing the optimization problem.
\cite{survey,landau}

\subsection{Preserve useful partial work without weakening publication}
The current complete-pass checkpoint is safe but coarse. A timeout late in a
large symbolic-host traversal can discard many independently certified local
improvements from that pass. Publishing a partially rebuilt tree safely requires
tracking the original context and already certified replacements, not simply
returning whatever local variable happened to be last assigned.

A persistent expression tree or a zipper with a replacement log would permit
finer checkpoints. A supervising process could store those checkpoints durably.
The revision retains complete-pass publication because it is easier to audit;
it does not disguise that decision as optimal recovery behavior.

\subsection{Optimize a small portfolio rather than a single presentation}
Strictly improving local moves can miss a path through a temporarily larger but
algebraically more useful expression. A bounded Pareto frontier over depth,
node count, coefficient size, and field degree could retain a few such states.
Only the final published representation would need to beat the original cost.
Every state would still need exact value certification, and the frontier would
need its own explicit size and work limits.

This is an opportunity for better search, not a reason to relax equality checks.
The same principle applies to fourth-root and root-of-unity presentations: the
allowed output grammar should be an explicit policy, especially when a user
wants real radicals only. The current grammar allows principal complex rational
powers and therefore measures a different search problem from real-only
square-root denesting.

\section{Conclusion}
The second corrected implementation has the right basic architecture for a
practical, incomplete denester: it treats discovery as fallible and places a
verification gate between a proposal and a replacement. The most valuable next
step is to apply that discipline consistently to the surrounding software
contracts.

The root recognizer must establish its claimed domain. Every search stage must
preserve the incumbent it receives. An index-reduced expression must be offered
even when recursion does nothing. A negative cache entry must not outlive the
resource conditions that created it. Limits must be checked where proposals
and expensive operations are actually generated, and reports must not perform
unbounded work after the budget has expired.

The accompanying revision implements those repairs while retaining the existing
search machinery. It also adds two mathematically transparent extensions:
bounded character reconstruction for small multiquadratic fields and an odd-index
trace--norm kernel in quadratic fields. Exact proofs, concrete witnesses,
reproducible source generation, independently executed checks, and an explicitly
unexecuted native suite make the proposal useful without overstating its
validation status.

\appendix
\section{Archive contents and rebuild procedure}
\begin{longtable}{@{}P{.40\textwidth}P{.53\textwidth}@{}}
\toprule
File & Purpose\\\midrule\endhead
\code{article.tex}, \code{article.pdf} & This self-contained article and its
rendered edition.\\
\code{StradFixed3.wl} & Standalone proposed package in a separate public context.\\
\code{StradFixed2.original.wl} & Exact reviewed baseline, retained for reproducibility.\\
\code{build_revision.py} & Hash-checking, anchor-checked source generator.\\
\code{revision_definitions.wl} & Replacement private definitions and new kernels.\\
\code{StradFixed2-to-3.diff} & Human-readable unified source comparison.\\
\code{StradFixed3Tests.wlt}, \code{run_tests.wls} & Native Wolfram tests and runner;
not executed in this review.\\
\code{independent_checks.py}, \code{independent_results.json} & Executed exact
mathematical/source checks and their full machine-readable results.\\
\code{README.md}, \code{source-manifest.json} & Usage, qualification notes,
source provenance, and relevant paths.\\
\code{SHA256SUMS.txt}, \code{build_pdf.sh} & Artifact-integrity hashes and PDF
rebuild commands.\\\bottomrule
\end{longtable}
To reproduce the proposed source and independent checks:
\begin{lstlisting}
python build_revision.py
python independent_checks.py
\end{lstlisting}
To rebuild this article, run pdfLaTeX twice from the archive directory. No
external images, bibliography database, local font files, or custom class are
required. The document uses standard TeX Live/MiKTeX packages and embedded
scalable Latin Modern fonts.

\section{Source locations for the principal findings}
All line references in this appendix refer to the archived, hash-verified
\code{StradFixed2.original.wl}, not to the generated revision.
\begin{longtable}{@{}P{.28\textwidth}P{.20\textwidth}P{.45\textwidth}@{}}
\toprule
Concern & Lines & Inspection target\\\midrule\endhead
Domain admission & 143--166 & Opaque objects admitted by precision/no-Real test.\\
Cost accounting & 168--181 & Subtractive radical count and atom bit-size scans.\\
Equality status & 183--221 & Numeric Different branch; exact-oracle fallback.\\
Acceptance/polishing & 224--255 & Single gate; size check after candidate entry.\\
Rationalization & 261--279 & Polynomial inverse and exact product certificate.\\
Multisurd ansatz & 345--375 & Sparse generator sets, FactorInteger, ansatz expansion.\\
Linear roots/index reduction & 420--455 & Together call, recursive-progress condition.\\
Multiplier expansion & 560--572 & Full discriminant factorization; batch guard placement.\\
Incumbent resets & 598--620 & Stage calls that omit the incumbent.\\
Memoization & 628--639 & Target-only key; unchanged-result storage.\\
Session/report & 692--728 & Scope of constrained evaluation and report cost scans.\\
Public dispatch & 736--742 & OptionsPattern tail versus validation contract.\\\bottomrule
\end{longtable}

\begin{thebibliography}{99}
\raggedright
\bibitem{source} Vladimir Reshetnikov, RadicalDenest repository,
\emph{StradFixed2.wl}, commit
\texttt{9000d6f533a68ee7659a6830ad25bba401f800a5}.
\url{https://github.com/VladimirReshetnikov/RadicalDenest/blob/9000d6f533a68ee7659a6830ad25bba401f800a5/src/corrected/StradFixed2.wl}.
The exact file is included in the archive.
\bibitem{prior} RadicalDenest repository, \emph{Unified analysis B},
\nolinkurl{src/code-review/unified-B/unified_analysis_B.tex}, same commit.
\url{https://github.com/VladimirReshetnikov/RadicalDenest/tree/9000d6f533a68ee7659a6830ad25bba401f800a5/src/code-review/unified-B}.
\bibitem{survey} RadicalDenest repository, \emph{Radical denesting: unified report},
\nolinkurl{docs/report/radical_denesting_unified.tex}, same commit.
\url{https://github.com/VladimirReshetnikov/RadicalDenest/tree/9000d6f533a68ee7659a6830ad25bba401f800a5/docs/report}.
\bibitem{editions} RadicalDenest repository, \emph{Corrected re-typesettings of six
articles}, \nolinkurl{docs/literature/RETYPESET_ARTICLES.md}, same commit.
\url{https://github.com/VladimirReshetnikov/RadicalDenest/blob/9000d6f533a68ee7659a6830ad25bba401f800a5/docs/literature/RETYPESET_ARTICLES.md}.
\bibitem{jeffrey} David J. Jeffrey and Albert D. Rich,
\emph{Simplifying Square Roots of Square Roots by Denesting}, chapter 4.
Consulted the repository's corrected 2026 re-typesetting, especially
Sections 4.4--4.5, at
\url{https://github.com/VladimirReshetnikov/RadicalDenest/blob/9000d6f533a68ee7659a6830ad25bba401f800a5/docs/literature/papers/jeffrey-Simplifying-square-roots-of-square-roots-by-denesting/retypeset-2026/denest.tex}.
\bibitem{bfht} Allan Borodin, Ronald Fagin, John E. Hopcroft, and Martin Tompa,
\emph{Decreasing the Nesting Depth of Expressions Involving Square Roots}
(1985). Consulted the repository's corrected 2026 re-typesetting, including
its explicit editorial qualifications:
\url{https://github.com/VladimirReshetnikov/RadicalDenest/blob/9000d6f533a68ee7659a6830ad25bba401f800a5/docs/literature/papers/bfht-Decreasing-the-nesting-depth-of-expressions-involving-square-root/retypeset-2026/symb85.tex}.
\bibitem{landau} Susan Landau, \emph{Simplification of Nested Radicals}, 1989
extended abstract. Consulted the corrected 2026 repository re-typesetting and
its editorial distinction from the final 1992 SIAM Journal on Computing paper,
DOI 10.1137/0221009:
\url{https://github.com/VladimirReshetnikov/RadicalDenest/blob/9000d6f533a68ee7659a6830ad25bba401f800a5/docs/literature/papers/landau92-Simplification-of-nested-radicals/retypeset-2026/landau1989.tex}.
\bibitem{wlroot} Wolfram Research, \emph{Root}, Wolfram Language documentation,
consulted September 2026. \url{https://reference.wolfram.com/language/ref/Root.html}.
\bibitem{wlzero} Wolfram Research, \emph{PossibleZeroQ}, particularly the
\code{Method -> "ExactAlgebraics"} contract, consulted September 2026.
\url{https://reference.wolfram.com/language/ref/PossibleZeroQ.html}.
\bibitem{wlreduce} Wolfram Research, \emph{RootReduce}, consulted September 2026.
\url{https://reference.wolfram.com/language/ref/RootReduce.html}.
\bibitem{wltime} Wolfram Research, \emph{TimeConstrained}, details on CPU accounting,
interrupt granularity, and abort protection, consulted September 2026.
\url{https://reference.wolfram.com/language/ref/TimeConstrained.html}.
\bibitem{wlmemory} Wolfram Research, \emph{MemoryConstrained}, consulted September
2026. \url{https://reference.wolfram.com/language/ref/MemoryConstrained.html}.
\bibitem{wlradicals} Wolfram Research, \emph{ToRadicals}, consulted September 2026.
\url{https://reference.wolfram.com/language/ref/ToRadicals.html}.
\end{thebibliography}
\end{document}
